% BHU PYQ -- Auto-generated & error-corrected
\documentclass[11pt,a4paper]{article}

\usepackage[a4paper,top=2.5cm,bottom=2.5cm,left=2.8cm,right=2.8cm,headheight=14pt]{geometry}
\usepackage{amsmath,amssymb}
\usepackage[T1]{fontenc}
\usepackage[utf8]{inputenc}
\usepackage{lmodern}
\usepackage{microtype}
\usepackage{fancyhdr}
\usepackage{enumitem}
\usepackage{hyperref}
\usepackage{lastpage}
\usepackage{parskip}

\hypersetup{colorlinks=true,urlcolor=black,linkcolor=black,
  pdftitle={BAS-101: Applied Statistics -- BHU PYQ},pdfauthor={Banaras Hindu University}}

\pagestyle{fancy}\fancyhf{}
\renewcommand{\headrulewidth}{0.4pt}\renewcommand{\footrulewidth}{0.4pt}
\fancyhead[L]{\small\textit{Banaras Hindu University}}
\fancyhead[R]{\small\textit{BAS-101: Applied Statistics}}
\fancyfoot[C]{\small Page \thepage\ of \pageref{LastPage}}

\newlist{parts}{enumerate}{2}
\setlist[parts,1]{label=(\alph*),leftmargin=*,itemsep=5pt,topsep=3pt}
\setlist[parts,2]{label=(\roman*),leftmargin=*,itemsep=3pt,topsep=2pt}
\newcommand{\pts}[1]{\hfill\textit{[\#1]}}

\begin{document}

\begin{center}
  {\large\bfseries BANARAS HINDU UNIVERSITY}\\[2pt]
  {\normalsize B.A. (Hons.) Semester I Examination, 2022-23}\\[6pt]
  \rule{0.8\linewidth}{0.5pt}\\[6pt]
  {\large\bfseries Paper No. BAS-101: Applied Statistics}\\[4pt]
  {\normalsize Time: 3 Hours\hfill Maximum Marks: 70}
\end{center}
\rule{\linewidth}{0.4pt}
\smallskip
\noindent\textit{Instructions: Answer \textbf{five} questions including Question~1 which is compulsory. Symbols have their usual meanings.}
\bigskip

BAS - 101 : Statistical Methods
N
\medskip\noindent\textbf{Question 1.}

\begin{parts}
  \item Write the definition of statistics with its limitations.
  \item Define the terms: Categorical Variable; Discrete and Continuous data; Interval and Ratio scale of measurement
\end{parts}
\medskip\rule{\linewidth}{0.2pt}

\medskip\noindent\textbf{Question 2.}

\begin{parts}
  \item What is.difference between Primary \& Secondary data?
  \item Discuss different ways of presenting the data with examples. a
\end{parts}
\medskip\rule{\linewidth}{0.2pt}

\medskip\noindent\textbf{Question 3.}

\begin{parts}
  \item Write different steps in preparing grouped frequency distribution table. ;:
  \item What are different ways of presenting the data graphically?
\end{parts}
\medskip\rule{\linewidth}{0.2pt}

\medskip\noindent\textbf{Question 4.}

\begin{parts}
  \item What is the purpose of measures of central tendency and variability?
  \item What are the different measures of central tendency?Discuss the measures with their merits and demerits.
\end{parts}
\medskip\rule{\linewidth}{0.2pt}

\medskip\noindent\textbf{Question 5.}

\begin{parts}
  \item Define coefficient of variation?What purpose does it serve?
  \item Discuss absolute and relative measures'of dispersion.
\end{parts}
\medskip\rule{\linewidth}{0.2pt}

\medskip\noindent\textbf{Question 6.}

\begin{parts}
  \item . Define moments. Distinguish between raw and central moments.
  \item What do you mean by positively and negatively skewed distribution. Also discuss various  measures of skewness.
\end{parts}
\medskip\rule{\linewidth}{0.2pt}

\medskip\noindent\textbf{Question 7.}

\begin{parts}
  \item What is scatter diagram?How is it useful in the study of correlation?
  \item Distinguish between linear and non-linear correlation. Mention important properties of correlation coefficient. Write short notes on any two of the following:
\end{parts}
\medskip\rule{\linewidth}{0.2pt}

\medskip\noindent\textbf{Question 8.}

\begin{parts}
  \item Measure of Kurtosis
  \item Yule's Coefficient of Association
  \item Uses of Regression Analysis 
\end{parts}

\vfill
\rule{\linewidth}{0.4pt}
\begin{center}\small
  \href{https://bhu-exam.vercel.app/}{Click here to visit our website}\\[4pt]\href{https://me-aryan.vercel.app/}{Click here to visit our contributor}\\[4pt]\href{https://quashan-me.vercel.app/}{Click here to visit our contributor}
\end{center}

\end{document}